\section{Limitations, Ethics, and Scope}

CiTT is educational design-assistance software. It is not patient-facing software and does not diagnose, monitor, treat, or control patient care. If a future version were used for patient-connected verification, clinical decision support, or device control, the scope and regulatory pathway would need reassessment.

The system has several technical limitations. Parsing can fail or misread diagrams. Agentic model generation can produce incomplete models. Students can overtrust simulation outputs. Simscape models may omit nonidealities that matter in real hardware or biology. Benchmark 2 retains symbolic values, so numeric outputs require additional parameters. Benchmark 3 uses educational scaled parameters and shows non-settling, saturation, and algebraic-loop warnings; these are limitations to teach from, not a claim of physiological fidelity.

The full Lab Delta CSV workflow remains pending without an external lab CSV. The no-tools Gemini baseline is a limited comparison artifact rather than an exhaustive study. Future work should add controlled user studies, broader benchmark suites, stronger model-check coverage, and instructor-facing review tools.

\section{Risk Mitigation}

CiTT mitigates overtrust by exposing structured specs, assumptions, limitations, model checks, focus/probe maps, and exported evidence. The UI should continue to make uncertainty visible. Instructors and students remain responsible for reviewing generated models and interpreting simulation evidence.
